\subsection*{The published comparison and its precise obstruction}

\statement{B.1}{Proved}{The closure-intersection step is an unproved implication}
The proofs of Proposition 2.4 of \citet{Boggi2026} and Theorem 5.1(iii) of
\citet{BoggiZalesskii2025} contain the same unproved passage. After the
per-element comparisons, the equality that remains to be established is
\[
 \bigcap_{x\in S}\overline{C_G(x)}
       =\overline{\bigcap_{x\in S}C_G(x)}
       =\overline{C_G(H)},\qquad H=\langle S\rangle .
\]
Here \(S\) is a finite generating set of the finite subgroup \(H\leq G\),
and all closures are in the specified profinite completion of \(G\).
The cited per-element comparisons alone give only the inclusion from right
to left. This observation does not refute either theorem.
The exact published scopes must be retained. In Proposition 2.4,
$G$ is assumed residually finite, good, and of finite virtual
cohomological type; the assertion is quantified over every finite
$H\leq G$. Its central conclusion is
$Z_{\widehat G}(\iota(H))=\overline{\iota(Z_G(H))}$.
In Theorem 5.1, a class $\mathscr C$ of finite groups is fixed.
The group is residually $\mathscr C$, $\mathscr C$-good, and of
finite virtual $p$-cohomological type at the relevant primes.
Every finite solvable $\mathscr C$-subgroup has a normalizer with
those same properties and with its full pro-$\mathscr C$ topology
induced from $G$. Part (iii) asserts
$C_{G_{\widehat{\mathscr C}}}(H)=\overline{C_G(H)}$
for every finite $\mathscr C$-subgroup $H$, while its normalizer
assertion is restricted to solvable $H$.
We audit these statements with their hypotheses, not a version
in which goodness alone is substituted for all of them.

\begin{proof}
For any subset \(S\) of any group, \(C(S)=\bigcap_{x\in S}C(x)\).
Consequently the left side of the displayed equality is already
\(C_{\widehat G}(H)\), once \(C_{\widehat G}(x)=\overline{C_G(x)}\)
has been established for \(x\in S\). Thus the missing reverse inclusion
is exactly the asserted finite-subgroup comparison. A closure operator
preserves inclusions and finite unions; it does not by its definition
commute with intersections. Neither proof supplies a separate argument
for that reverse inclusion at this point.

We checked Proposition 2.4 also in the published Springer text, not only
in the supplied arXiv source. Its displayed chain is unchanged at this
step. In the supplied source of Boggi--Zalesskii, the corresponding chain
is in the last paragraph of the proof of Theorem 5.1. Their proof of the
normalizer assertion for solvable subgroups invokes Lemma 2.3 and
Theorem 5.1(ii) separately and does not invoke this chain.
\end{proof}

\statement{B.2}{Proved}{The normalizer comparison repairs the solvable centralizer case}
Let \(G\) be residually finite and \(H\leq G\) finite. If
\(N_{\widehat G}(H)=\overline{N_G(H)}\), then
\[
 C_{\widehat G}(H)=\overline{C_G(H)}.
\]
In particular the centralizer conclusion of Boggi--Zalesskii,
Theorem 5.1, is justified for solvable \(H\) under that theorem's two
hypotheses, by its independently proved normalizer assertion.
\begin{proof}
Conjugation defines a continuous homomorphism
\(N_{\widehat G}(H)\to\Aut(H)\), where the target is finite and discrete.
Its kernel is \(C_{\widehat G}(H)\), hence is open in the normalizer.
For a dense subgroup \(D\) of a topological group \(L\), and an open
subgroup \(V\leq L\), the intersection \(D\cap V\) is dense in \(V\):
every open subset of \(V\) is open in \(L\). Apply this with
\(D=N_G(H)\), \(L=N_{\widehat G}(H)\), \(V=C_{\widehat G}(H)\).
Then \(D\cap V=C_G(H)\), proving the claim.

This repair retains the normalizer hypotheses. Proposition 2.4 in
Boggi's paper has fewer stated hypotheses, so a citation of the solvable
case of Theorem 5.1 without verifying its second hypothesis would not
repair Proposition 2.4 in that generality.
\end{proof}

\statement{B.3}{Proved}{The exact coset obstruction}
Let \(U\trianglelefteq G\) have finite index and \(U\cap H=1\). Write
\(\widehat U=\overline U\leq\widehat G\). For
\(\bar g\in C_{G/U}(\bar H)\), choose \(g\in G\) mapping to \(\bar g\) and put
\[
 a_g(h)=ghg^{-1}h^{-1}\in U\qquad(h\in H).
\]
With \(H\) acting on \(U\) by conjugation, this is a nonabelian
one-cocycle. The class \([a_g]\in H^1(H,U)\) is independent of the
choice of \(g\) in the left coset \(Ug\). Then
\[
\begin{split}
 \bar g\in\im\bigl(C_G(H)\to G/U\bigr)
 &\ \Longleftrightarrow\ [a_g]=1\text{ in }H^1(H,U),\\
 \bar g\in\im\bigl(C_{\widehat G}(H)\to G/U\bigr)
 &\ \Longleftrightarrow\ [a_g]=1\text{ in }H^1(H,\widehat U).
\end{split}
\]
Hence centralizer equality is equivalent to triviality of the completion
kernel on these \emph{coset obstruction classes}, at every \(U\).
\begin{proof}
Direct multiplication gives
\(a_g(hk)=a_g(h)\,h(a_g(k))\). For \(u\in U\), one has
\(a_{ug}(h)=u\,a_g(h)\,h(u)^{-1}\), the usual change by a coboundary.
Moreover \(ug\) centralizes \(H\) precisely when
\(a_g(h)=u^{-1}h(u)\) for all \(h\). The same calculation in
\(\widehat U\) proves the second equivalence. Finally a closed subgroup
of a profinite group is determined by its images in all finite
quotients. This proves the last assertion.

All classes in \(H^1(H,U)\) classify marked complements in \(U\rtimes H\).
Only the classes displayed above are automatically obtained by conjugating
the given marked inclusion by an element of \(G\). Thus the unrestricted
\emph{every complement} formulation in R5 requires an additional
conjugacy/marking comparison. It is a sufficient condition, but it is
not obtained merely by rewriting the displayed centralizer equality.
The subgroup-conjugacy criteria in Boggi--Zalesskii, Lemmas 2.2--2.3,
likewise retain a separate conjugacy-distinguished hypothesis.
\end{proof}

\subsection*{The icosahedral quotient and its outer centralizer}

\statement{B.4}{Proved}{The level-five kernel is characteristic}
Every continuous epimorphism
\(\widehat B=\widehat{C_2*C_3}\twoheadrightarrow A_5\) has the same
kernel \(\widehat K\). In particular \(\widehat K\) is characteristic
in \(\widehat B\).
\begin{proof}
An epimorphism sends the two free factors to elements of orders
two and three. In \(A_5\) their product must have order five:
if its order were two or three, the image would be a quotient of
the spherical triangle group of type \((2,3,2)\) or \((2,3,3)\),
of orders six and twelve respectively. Thus the epimorphism factors
through \(\langle s,r\mid s^2=r^3=(sr)^5=1\rangle=A_5\).
Since this quotient already has order sixty, the induced surjection
from it to \(A_5\) is an automorphism. Its kernel is therefore
independent of the epimorphism.

The program \texttt{code/check\_icosahedral\_groups.py} enumerates all
300 ordered pairs of the prescribed orders. The counts, by
(generated subgroup order, product order), are
\((6,2):60\), \((12,3):120\), \((60,5):120\).
An independent Todd--Coxeter computation gives order sixty for
the triangle presentation; its complete coset table is included.
These are checks of the finite argument, not an inference from a
noncofinal collection of quotients.
\end{proof}

\statement{B.5}{Proved}{A centerless extension comparison}
Let \(1\to L\to E\to H\to1\) be an exact sequence of profinite groups,
with \(L\) centerless, \(C_E(L)=1\), \(H\) centerless, and \(L\)
characteristic in \(E\). View \(H\) as a subgroup of \(\Out(L)\)
through conjugation. Then restriction induces a canonical isomorphism
\[
 \ker\bigl(\Out(E)\longrightarrow\Out(H)\bigr)
       \ \simeq\ C_{\Out(L)}(H).
\]
The groups carry their profinite automorphism topologies when they are
topologically finitely generated.
\begin{proof}
Conjugation identifies \(E\) with the inverse image of \(H\) under
\(\Aut(L)\to\Out(L)\). Indeed it identifies its kernel with
\(\Inn(L)\), and both quotients are the given \(H\).
If \([\alpha]\in C_{\Out(L)}(H)\), conjugation by any representative
\(\alpha\in\Aut(L)\) preserves this inverse image, hence gives a
continuous automorphism of \(E\) inducing the identity on \(H\).
Changing \(\alpha\) by an inner automorphism of \(L\) changes the
extension automorphism by the corresponding inner automorphism
of \(E\).

Conversely, change a representative of a class in the displayed
kernel by an inner automorphism of \(E\) so that its action on \(H\)
is the identity. Its restriction to \(L\) then centralizes \(H\)
outerly. Two such adjusted representatives differ by conjugation
by an element whose image is in \(Z(H)=1\), and hence differ by an
inner automorphism of \(L\). The two constructions are inverse.
They are induced by continuous restriction and conjugation; in the
finitely generated case a continuous bijection of the resulting
compact Hausdorff groups is a homeomorphism.
\end{proof}

\statement{B.6}{Proved}{The peripheral icosahedral comparison}
Let \(\Out_{\partial}(\widehat K)\) denote continuous outer
automorphisms permuting the twelve conjugacy classes of cusp inertia,
with a common unit multiplier. There is a canonical isomorphism
\[
 C_{\Out_{\partial}(\widehat K)}(A_5)
 \ \simeq\
 \{\Phi\in\GTzero:\lambda(\Phi)\equiv1\text{ or }-1\pmod5\}.
\]
The isomorphism is between outer automorphism groups; a chosen
automorphism representative is still ambiguous by inner conjugation.
\begin{proof}
The open subgroup \(\widehat K\) is free profinite of rank eleven and
has trivial center. The centralizer in \(\widehat B\) of this open
subgroup is trivial: an element centralizing it has a power in its
center and hence is torsion, whereas a nontrivial torsion element
of a free profinite product has finite centralizer in its finite
free factor. Thus B.5 applies.

For a cusp subgroup \(I=\overline{\langle\tau^5\rangle}\leq\widehat K\),
its normalizer in \(\widehat K\) is \(I\). This is the
peripheral normalizer assertion for a primitive free generator,
also obtainable by writing the free profinite group as
\(I\amalg\widehat F_{10}\). The action of
\(A_5=B/K\) on the twelve cusp classes has stabilizer \(C_5\).
It follows that
\[
 N_{\widehat B}(I)=\overline{\langle\tau\rangle}.
\]
Indeed any element of the left side has image in that cusp
stabilizer, so can be multiplied by a power of \(\tau\) to lie
in \(\widehat K\); the remaining element is in \(I\).
Consequently the extension construction in B.5 preserves the
twelve cusp classes precisely when its automorphism of
\(\widehat B\) preserves the conjugacy class of
\(\overline{\langle\tau\rangle}\).

The cusp-preserving outer group of \(\widehat B\) identifies with
\(\GTzero\). The comparison is restriction to the characteristic
level-two free subgroup, followed by its three-cusp identification:
the quotient \(S_3\) is centerless, every automorphism of \(S_3\) is
inner, and B.5 identifies this outer group with the centralizer of
the geometric \(S_3\) in the three-cusp outer group.
Harbater--Schneps, Main Theorem(a), with Propositions 3--4 and
Theorem 2, identifies that centralizer with \(\GTzero\);
the cusp comparison follows by the same normalizer argument with
\(\tau^2\).

Finally \(\Aut(A_5)=S_5\), and its outer involution exchanges the
two conjugacy classes of elements of order five. A generator \(t\)
of a Sylow five subgroup is conjugate in \(A_5\) to \(t^\lambda\)
exactly for \(\lambda\equiv\pm1\pmod5\).
The induced automorphism of \(A_5\) is therefore inner exactly for
these multipliers. Apply B.5.
\end{proof}

\statement{B.7}{Proved}{Purity and the discrete centralizer}
The multiplier-one subgroup in B.6 acts trivially on the twelve
puncture labels. Moreover
\[
 C_{\Gamma_{0,[12]}}(A_5)=1.
\]
Under the faithful genus-zero profinite Dehn--Nielsen representation,
the full profinite centralizer identifies exactly with the group
\(\mathcal K\) of the question, with no nontrivial finite discrete factor.
\begin{proof}
Choose the extension automorphism of \(\widehat B\) inducing the
identity on \(A_5\). Write
\(\Phi(\tau)=b\tau^\lambda b^{-1}\). In \(A_5\),
\(t=\bar b\,t^\lambda\bar b^{-1}\). If
\(\lambda\equiv1\pmod5\), then
\(\bar b\in C_{A_5}(t)=\langle t\rangle\). Thus
\(b\in\widehat K\,\overline{\langle\tau\rangle}\), and the cusp
class of \(\tau^5\) is fixed in \(\widehat K\).
Equivariance and transitivity fix every cusp. This applies in
particular when \(\lambda=1\).

For discrete orientation-preserving mapping classes, the normalizer
comparison for the regular branched covering identifies the
normalizer modulo its deck group with the mapping classes of the
quotient preserving the covering. The quotient has exactly the
three distinguished points of orders \(2,3,\infty\). Their orders
are different, so they are fixed individually, and the pure
three-point mapping class group is trivial. The discrete normalizer
is therefore \(A_5\), and its centralizer is \(Z(A_5)=1\).
This argument uses the discrete covering theorem, not the
profinite centralizer comparison under audit.

The geometric mapping class completion acts with multiplier one.
Its faithful embedding into \(\Out_{\partial}(\widehat K)\)
(the genus-zero congruence theorem) and B.6 now identify its
centralizer with the stated \(\mathcal K\).
R1 makes \(\mathcal H\), hence \(\mathcal K\), torsion-free.
Thus in this case a finite centralizer would be a trivial centralizer.
\end{proof}
